\chapter{Выводы}
\label{ch:conclusions}

    %Вывод - результат аналитической деятельности, а не констатация факта. Другими словами, чем вывод в анотации, более развёрнуто.
    
    Измерен коэффициент теплопроводности воздуха при атмосферном давлении в зависимости от температуры в диапазоне от $33$ до $80 \celsius$ с помощью термостата, заполненного воздухом, и источника тепла. Отмечено, что с ростом температуры $T$ коэффициент теплопроводности имел степенную зависимость $T^{\beta}$ и изменялся от $26.2 \cdot 10^{-3}$ до $29.7 \cdot 10^{-3} \frac{\text{Вт}}{\text{м} \cdot \text{К}}$, при этом $\beta = 0.87 \pm 0.03$. Полученное значение $\beta$ не соответствует $\beta = 0.5$, которое соответствует простейшей модели твёрдых шариков. Это показывает, что использование простейшей модели твёрдых шариков при оценке коэффициента теплопроводности даёт большую погрешность. Это связано с тем, что в модели твёрдых шариков эффективное сечение столкновений молекул считается постоянным, а на практике является медленно убывающей функцией от T. В работе приведена методика, позволяющая более точно определить, как изменяется коэффициент теплопроводности в зависимости от температуры, экспериментально.